\head{L1 observations}
The historical standalone threshold consumes 32-bit squared magnitudes; the current fused stage consumes 16-bit L1 magnitudes. The separate 64-pixel threshold loop needs eight port-5 pack/permutation operations, giving an ideal eight pixels/cycle. Counting four comparison/classification operations per pixel gives 32 useful lane-ops/cycle on that convention. At width 512, unrolling to 128 outputs reached 0.215 cycles/pixel (18.6 useful ops/cycle), versus 0.309 for 64. Production thresholds stay inside NMS, whose output contract and instruction mix differ.
\head{Independent scan versus dependent queue}
For the exact 32-pixel scan, 21 modeled load-port service units and two load ports give $320/(21/2)=30.5$ byte lane-ops/cycle under the stated split assumption. It reached 17.29 (57\%). The exact scan touched 6,740 bytes; ports 2 and 3 had the highest dispatch rates (0.459 and 0.443 per cycle). Pairing two scan rows improved the seed-free exact-scan median from 0.578 to 0.515 cycles/pixel. On a $258\times8$ slab, pairing improved sparse tracking from 3.496 to 2.753, but regressed dense tracking from 54.679 to 60.086. The queues contained 18 versus 1129 entries. The older 4-row seed predictor has a load service floor of $27/(2\times128)=0.1055$ cycles/pixel; it has a different output contract. A fixed per-pixel arithmetic ceiling is inappropriate for complete queue traversal.
